\documentclass[11pt]{article}
\usepackage[margin=1in]{geometry}
\usepackage{amsmath,amssymb,amsthm}
\usepackage{booktabs}
\usepackage{array}
\usepackage{float}
\usepackage{hyperref}
\usepackage{xcolor}

\newtheorem{theorem}{Theorem}
\newtheorem{proposition}{Proposition}
\newcommand{\systemu}{ELARA-U}
\newcommand{\stack}{\textsc{Stack}}
\newcommand{\auto}{\textsc{AutoSelect}}
\newcommand{\relg}{\textsc{RelGate}}
\newcolumntype{L}[1]{>{\raggedright\arraybackslash}p{#1}}

\title{\textbf{ELARA-U: A Strong Stacking Baseline, a Rigorous Negative Result,\\
and a Mechanism for Reliability-Aware Cross-Domain Anomaly Detection}\\[4pt]
\large Thesis Chapter (verified evidence edition)}
\author{}
\date{}

\begin{document}
\maketitle

\begin{abstract}
This chapter studies whether \emph{reliability-aware meta-routing}---selecting,
fusing, or abstaining among anomaly-detection experts using deployment-time
reliability signals---improves universal anomaly detection over simple baselines. On
a broad benchmark of \textbf{123 tasks across five families} (tabular, image-OOD,
text, cyber, fraud) we establish three verified results. (1) A parameter-light
rank-normalized logistic \emph{stack} is a strong, under-reported universal baseline:
it beats per-task validation selection by $+0.036$ AUROC (CI $[0.023,0.050]$) and the
best fixed detector by $+0.075$ ($[0.052,0.101]$), and exceeds the best-single-detector
oracle. (2) Reliability-aware routing does \emph{not} beat it: across three
degradation regimes and two independent router designs, with clean ablations,
reliability features carry no significant signal beyond validation quality and
detector identity, and an explicit reliability gate \emph{hurts} ($-0.037$,
$[-0.050,-0.026]$). (3) A controlled simulation shows reliability gating helps
\emph{iff} a channel can fail independently while others stay clean; single-modality
detectors share one input and violate this, which explains (2) and motivates a
\emph{pre-registered} multimodal test for the regime that satisfies it. The
contribution is an honest map of where reliability-aware anomaly routing helps, where
it does not, and why---not a claim of universal state of the art.
\end{abstract}

\section{Executive Summary}
No single anomaly detector dominates across domains; this drives \emph{negative
transfer} when a detector chosen on one domain is deployed on another. The obvious
remedy is to choose per task. We show that the strongest simple remedy is not a
reliability-aware router but a plain rank-normalized stack, and that the popular
reliability machinery does not improve on it on real single-modality data. We then
identify the exact structural condition under which reliability gating \emph{can}
help and pre-register a test for it.

\paragraph{Evidence level (honest).}
\begin{table}[H]\centering\small
\caption{What is established versus pending. All ``established'' rows are verified,
leakage-free, with paired-bootstrap confidence intervals over 123 tasks.}
\begin{tabular}{L{0.40\linewidth}L{0.32\linewidth}L{0.20\linewidth}}
\toprule
\textbf{Finding} & \textbf{Basis} & \textbf{Status} \\
\midrule
Stack $>$ validation selection $>$ fixed detector & 123 tasks, 5 families, clean splits & Established \\
Reliability features add no gain (3 regimes, 2 designs) & Clean ablations, bootstrap CIs & Established \\
Mechanism: gain needs independent channel failure & Controlled simulation & Established (simulation) \\
Reliability helps on real multimodal data & MVTec-3D / Real-IAD, GPU & Pre-registered, pending \\
Universal sealed-holdout SOTA claim & New external families & Not claimed \\
\bottomrule
\end{tabular}
\end{table}

\section{Research Questions and Contributions}
\textbf{RQ1.} Can a cross-domain decision layer beat the best fixed detector family
over many anomaly tasks? \textbf{RQ2.} Can it beat per-task validation-AUROC
selection? \textbf{RQ3.} Do deployment-time \emph{reliability} features (drift,
disagreement, dispersion, calibration) add value over validation quality and detector
identity? \textbf{RQ4.} If not, why, and under what condition would they?

\paragraph{Contributions.} (i) A strong, reproducible stacking baseline for universal
anomaly detection and the finding that it beats validation selection (RQ1, RQ2,
affirmative). (ii) A rigorous negative answer to RQ3 with two router designs and three
regimes, including a cautionary analysis of a parallel pipeline whose ``ablation''
was vacuous. (iii) A mechanistic answer to RQ4: reliability gating requires
independent channel failure, demonstrated in simulation and consistent with two
operating-boundary theorems. (iv) A pre-registered, falsifiable multimodal experiment
targeting the one regime that satisfies the condition.

\section{Problem Definition and Metrics}
For task $t$, detectors $D_1,\dots,D_M$ produce scores $s_m(x)$. Each task is split
into train (fit detectors, unlabeled), validation (labeled), and test (labeled,
held out), stratified $50/25/25$, standardized on train. $A_V(m),A_T(m)$ are
validation/test AUROC; $m^\star=\arg\max_m A_T(m)$ is the best-single-detector oracle
(a \emph{selection} ceiling---fusion may exceed it). Regret is $A_T(m^\star)-A_T(h)$.
We report average rank, mean AUROC, regret, ECE, per-family rank, and negative-transfer
rate, all with paired bootstrap CIs over tasks ($10^4$ resamples).

\paragraph{Leakage discipline.} No test labels fit any router, stacker, gate, or
threshold. Validation labels are permitted (available at deploy time). Unlabeled
test-score \emph{distributions} may inform reliability features---deployment
monitoring, marked as such, not label use. We state explicitly: the stack is
\emph{not} label-free; it uses validation labels.

\section{From ELARA to \systemu{}}
The predecessor system (ELARA / Reliability-Gated Attention) fused RGB and depth for
industrial anomaly detection using a validation-derived reliability gate. Its
transfer claim did not generalize and was honestly downgraded. \systemu{} re-poses the
question at benchmark scale and as a \emph{decision layer over a detector zoo}: rather
than asserting a reliability gate works, it \emph{tests} whether reliability signals
help selection/fusion across many tasks. This chapter's negative result is the
benchmark-scale resolution of the predecessor's open question, and its mechanism
section explains why the predecessor's multimodal setting is the only place the gate
could have helped.

\section{Methodology: the meta-router and its reliability variants}
We compare each \textbf{fixed} detector; \auto{} (per-task $\arg\max_m A_V(m)$);
\textbf{rank\_mean} and \textbf{cw\_mean} (naive ensembles); \stack{}, a logistic
regression on \emph{rank-normalized validation scores} $r_m(x)=\mathrm{rank}_m(x)/n$
applied to rank-normalized test scores; and reliability variants. The \textbf{reliability
gate} $\relg{}$ reweights channels by validation quality times an unlabeled test-time
consensus-agreement term; its \textbf{ablation} uses validation quality only. A
\textbf{learned router} (gradient-boosted regressor over detector identity, $A_V$,
dataset meta-features, and---in the reliability arm---score dispersion, gini,
skew/kurtosis, inter-detector disagreement, and val$\to$test KS drift) predicts
per-detector test AUROC under leave-task-out and leave-family-out cross-validation.

\paragraph{Detector zoo.} ECOD, COPOD, IForest, LOF, KNN, OCSVM (pyod), fit per task on
the unlabeled train split. Reliability/router features are computed from validation
scores and unlabeled test-score distributions only.

\section{Theory Stack}\label{sec:theory}
The theory characterizes the boundary rather than asserting universality (proofs in
Appendix~\ref{sec:proofs}).

\begin{theorem}[Clean ceiling]\label{thm:clean}
If validation and test are i.i.d.\ and detectors are calibrated on validation, no
reliability reweighting that depends only on the label-free test marginal improves
expected AUROC over the validation-optimal combination. Validation quality is a
sufficient statistic; reliability features are redundant.
\end{theorem}

\begin{theorem}[Independent-failure gain]\label{thm:stress}
If at deployment a strict subset of channels loses discriminative signal while the
complement ($\ge 2$ channels) stays informative, and the failure is detectable from
unlabeled test-score disagreement, then consensus-agreement reweighting attains
strictly lower expected regret than both equal-weight fusion and stale validation
selection.
\end{theorem}

Informally, \systemu{} helps only when stale-validation regret $\Delta$ exceeds the
router's prediction error $2\eta$; on i.i.d.\ near-ceiling data $\Delta\approx 0$ and
the router cannot help (Theorem~\ref{thm:clean}). It can help when deployment changes
the ranking in a \emph{detectable, partial} way (Theorem~\ref{thm:stress}).

\section{Benchmark Suite and Evidence Contract}
123 tasks: tabular (43, ADBench Classical), image-OOD (61, ADBench CV/ResNet18), text
(13, ADBench NLP/BERT), cyber (5, UNSW-NB15 splits), fraud (1). The evidence contract
separates: (a) average-rank improvement over the best fixed detector; (b) improvement
over \auto{}; (c) regret-to-oracle; (d) negative-transfer rate; (e) calibration; (f)
leave-family-out; (g) a no-test-label routing audit. A claim passes only with a
positive bootstrap CI on the relevant contrast.

\section{Experimental Evaluation}
\subsection{Headline: a strong stacking baseline (verified, clean splits)}
Table~\ref{tab:main} reports the main comparison on the clean benchmark---\emph{no
injected shift}, because the previously-used constant-offset degradation is
AUROC-invariant and therefore does not make validation stale (see
\S\ref{sec:validity}). \stack{} is the best non-oracle strategy and exceeds the
best-single-detector oracle ($0.784$ vs $0.756$): fusion beats selection.

\begin{table}[H]\centering
\caption{Main benchmark, 123 tasks / 5 families. Average rank among 11 non-oracle
strategies (lower better). Leakage-free.}
\label{tab:main}
\begin{tabular}{lrrrr}
\toprule
Strategy & Avg rank & Mean AUROC & Regret & ECE \\
\midrule
\stack{} (rank-norm logistic) & \textbf{2.53} & \textbf{0.784} & --- & 0.271 \\
\auto{} (val-AUROC) & 3.38 & 0.748 & 0.008 & 0.393 \\
\stack{}+reliability gate & 4.00 & 0.746 & 0.010 & 0.325 \\
fixed/KNN (best fixed) & 5.42 & 0.709 & 0.048 & --- \\
rank\_mean & 6.53 & 0.707 & 0.049 & 0.407 \\
cw\_mean & 6.66 & 0.708 & 0.048 & 0.396 \\
fixed/\{LOF,OCSVM,IForest,ECOD,COPOD\} & 6.4--8.7 & 0.68--0.70 & 0.05--0.08 & --- \\
\midrule
oracle (best single detector) & --- & 0.756 & 0.000 & --- \\
\bottomrule
\end{tabular}
\end{table}

\begin{table}[H]\centering
\caption{Key contrasts (paired bootstrap over 123 tasks).}
\label{tab:contrasts}
\begin{tabular}{lrrc}
\toprule
Contrast & $\Delta$AUROC & 95\% CI & Significant \\
\midrule
\stack{} $-$ \auto{}            & $+0.036$ & $[0.023,\,0.050]$ & yes \\
\stack{} $-$ best fixed         & $+0.075$ & $[0.052,\,0.101]$ & yes \\
\auto{} $-$ best fixed          & $+0.040$ & $[0.023,\,0.061]$ & yes \\
\stack{}+gate $-$ \stack{} (reliability ablation) & $-0.037$ & $[-0.050,\,-0.026]$ & \textbf{hurts} \\
\bottomrule
\end{tabular}
\end{table}

Per-family average rank shows the stack is robust, not driven by one family: tabular
$2.41$, image-OOD $2.49$, cyber $1.00$, text $3.69$; it loses only on fraud ($n{=}1$,
underpowered, excluded from family conclusions). This answers RQ1 and RQ2
affirmatively---and the gate already hints at RQ3's negative.

\subsection{The reliability negative (RQ3)}\label{sec:negative}
Table~\ref{tab:negative} tests the reliability premise with two router designs and
three regimes, each with a clean reliability-features-removed ablation. A learned
selector with the full reliability battery is \emph{worse} than \auto{} on i.i.d.\
data, and the ablation is null. Under uniform deployment shift a learned router beats
\emph{stale} selection, but leave-one-out attribution shows the gain is from
\emph{detector-identity priors} (which architectures survive corruption), not
reliability features (ablation null at every severity). Under heterogeneous per-task
missingness---the regime most favorable to reliability---the ablation flickers around
zero non-monotonically and never reaches significance.

\begin{table}[H]\centering
\caption{Reliability-feature ablations. $\Delta=$ (with reliability) $-$ (without);
$>0$ means reliability helps. None significant; the stack gate is significantly
negative.}
\label{tab:negative}
\begin{tabular}{llrc}
\toprule
Regime / design & comparison & $\Delta$AUROC & CI excludes 0? \\
\midrule
i.i.d., learned selector & rel.\ vs no-rel.\ & $+0.001$ & no \\
i.i.d., learned selector & rel.\ vs \auto{} & $-0.005$ & yes (worse) \\
uniform shift (sev.\ 3), learned & rel.\ vs no-rel.\ & $+0.002$ & no \\
heterog.\ missingness (frac.\ 0.7), learned & rel.\ vs no-rel.\ & $+0.015$ & no \\
123-task benchmark, stack gate & gate vs \stack{} & $-0.037$ & yes (worse) \\
\bottomrule
\end{tabular}
\end{table}

\paragraph{A cautionary note on ablations.} A parallel pipeline reported a polished
``reliability ablation'' in which every variant (full, no-reliability, no-drift,
no-disagreement, no-calibration) returned byte-identical metrics, because all were
silently mapped to the same model---the reliability features were never in the method.
An ablation that does not ablate is, at a glance, indistinguishable from one showing
``all components contribute.'' We surface it as a methodological hazard.

\subsection{Mechanism: reliability needs an independent clean channel (RQ4)}\label{sec:mechanism}
We isolate the cause in a controlled simulation (Table~\ref{tab:mechanism}) with two
regimes differing only in how degradation acts: \textbf{SHARED} weakens all channels
together (the single-modality analogue), \textbf{INDEPENDENT} fails one channel while
the rest stay clean (genuine multimodal). Methods are identical across regimes.

\begin{table}[H]\centering
\caption{Mechanism simulation (300 trials/regime, paired bootstrap).}
\label{tab:mechanism}
\begin{tabular}{lrrr}
\toprule
 & \relg{} $-$ mean & \relg{} $-$ \auto{} & \relg{} $-$ ablation \\
\midrule
SHARED (all weaken together) & $-0.016$ (ns) & $+0.081$ & $-0.019$ (ns) \\
INDEPENDENT (one fails, rest clean) & $+0.026$ \checkmark & $+0.144$ \checkmark & $+0.024$ \checkmark \\
\bottomrule
\end{tabular}
\end{table}

Under SHARED degradation reliability gating cannot help (no clean channel), matching
the single-modality benchmark. Under INDEPENDENT failure it beats equal-weight fusion,
stale selection, and its own no-test-time ablation, all CI$>0$: the unlabeled
test-time reliability signal, not validation quality, carries the gain. The negative
result is thus a statement that single-modality detectors violate the enabling
condition of Theorem~\ref{thm:stress}, not a refutation of reliability gating in
general.

\section{Validity and Leakage Audit}\label{sec:validity}
\begin{table}[H]\centering\small
\caption{Validity audit (honest controls and remaining work).}
\begin{tabular}{L{0.26\linewidth}L{0.44\linewidth}L{0.20\linewidth}}
\toprule
\textbf{Risk} & \textbf{Control} & \textbf{Remaining} \\
\midrule
Test-label routing & Routing uses validation labels + unlabeled test scores only; test labels score after routing. & Automated no-leakage assertions on every router path \\
Inert ``shift'' & The constant-offset degradation is AUROC-invariant; the headline now uses the clean benchmark. & Use natural shift only \\
Synthetic mechanism & Independent-vs-shared result is a controlled simulation, not real data. & The pre-registered multimodal test \\
Calibration & ECE reported; not a fitted probability calibrator. & Add isotonic + Brier/NLL \\
Family coverage & Leave-family-out used; only 5 families; fraud $n{=}1$. & More families; exclude $n{=}1$ from family claims \\
Ablation honesty & Reliability ablations decisive and \emph{negative}; gate hurts. & Claim is narrowed to rank stacking (correct) \\
\bottomrule
\end{tabular}
\end{table}

\section{Pre-registered Multimodal Test}\label{sec:multimodal}
Theorem~\ref{thm:stress} and \S\ref{sec:mechanism} make a falsifiable prediction for
RGB$+$3D anomaly detection, where a modality can fail independently. We pre-register
(frozen decision rule, no post-hoc retuning): under independent per-modality
degradation, reliability-gated fusion must beat (H1) equal-weight fusion, (H2) stale
validation selection, and (H3) its no-test-time ablation, each with paired-bootstrap
CI$>0$ over (category $\times$ seed) units. A CPU positive control with the correct
design (two modalities, two detectors each, modality-level failure) passes decisively:
H1 $+0.094$, H2 $+0.264$, H3 $+0.097$, all CI$>0$. The real test runs on MVTec-3D /
Real-IAD. A pass yields a scoped novel-method claim; a failure is reported as the final
boundary, with no novel claim and no retuning to chase a pass.

\section{Honest Scope, Limitations, and Path Forward}
The positive baseline and the reliability negative hold for single-modality
feature-vector tasks only; no industrial-vision or time-series claim is made. The
reliability gate tested is one (consensus-agreement) instantiation; three regimes and
two designs make a different parameterization helping on this data unlikely but not
impossible. The mechanism is a simulation; it motivates, not proves, the multimodal
claim---hence pre-registration. Path to a stronger result: (i) run the multimodal test;
(ii) add natural-shift and time-series families; (iii) sealed external holdouts; (iv)
fitted calibration with Brier/NLL.

\section{Reproducibility}
One score archive (123 tasks $\times$ 6 detectors) persists all val/test scores; every
result re-derives from it without re-running detectors. Scripts: \texttt{honest\_benchmark.py}
(Table~\ref{tab:main},\ref{tab:contrasts}), \texttt{learned\_router\_ablation.py},
\texttt{shift\_stress\_ablation.py}, \texttt{heterogeneous\_degradation\_ablation.py}
(Table~\ref{tab:negative}), \texttt{synthetic\_multimodal\_poc.py}
(Table~\ref{tab:mechanism}), \texttt{multimodal\_reliability\_experiment.py} +
\texttt{MULTIMODAL\_RELIABILITY\_PROTOCOL\_v1.md} (\S\ref{sec:multimodal}). Fixed seeds;
deterministic missingness via stable hashing; no test labels in any routing path.

\section{Conclusion}
On a broad anomaly benchmark, a simple rank-normalized stack is a strong, under-reported
universal baseline, and reliability-aware routing does not improve on it. We trace this
to one structural condition---reliability gating requires an independently failing
channel, which single-modality detection cannot provide---reconcile it with the
operating-boundary theory, and pre-register the multimodal test for the regime that
satisfies the condition. The result is an honest boundary map, not a universal SOTA
claim.

\appendix
\section{Proofs}\label{sec:proofs}
\begin{proof}[Proof sketch of Theorem~\ref{thm:clean}]
With i.i.d.\ splits the conditionals $p(s\mid y)$ are identical across splits. The
AUROC of any scalarization $w^\top r(x)$ is a functional of these conditionals, so the
validation-AUROC-optimal $w$ maximizes expected test AUROC. A reweighting depending
only on the label-free marginal $p(s)$ carries no information about $p(s\mid y)$ beyond
what validation used, hence cannot improve $w$ in expectation. $\square$
\end{proof}
\begin{proof}[Proof sketch of Theorem~\ref{thm:stress}]
Partition channels into informative $\mathcal{I}$ ($|\mathcal{I}|\ge2$) and failed
$\mathcal{F}$. Failed channels are independent of $y$, so they add zero-mean noise to
equal-weight fusion (raising regret) and exhibit low rank-correlation with the
$\mathcal{I}$-consensus. Weights $w_m\propto q_m a_m$ with agreement $a_m$ drive
$w_{\mathcal F}\to0$, recovering fusion on $\mathcal I$, whose expected AUROC strictly
exceeds both the diluted equal-weight fusion and any single (possibly stale-selected)
channel. $\square$
\end{proof}

\end{document}
